\section{The finite graph reduction}
\label{sec:construction}

We first construct a graph from an arbitrary Label Cover instance and
fixed finite parameters. A satisfying labeling will give its
three-coloring. The following sections analyze independent sets and
determine which parameter choices make their density small.

Fix integers $r\ge s\ge2$ and $m\ge1$, an even positive integer $P$,
a rational $\lambda\in(0,1)$, and a rational probability distribution
$\mu$ on $\binom{[r]}s$. Set $D=mP$. The coefficient values and their
probabilities are
\begin{equation}\label{eq:coefficient-law}
 \mathcal T_m=\{0,1/m,\ldots,1\},\qquad
 p_k=\Pr(t=k/m)=\frac{\lambda^k}{\sum_{\ell=0}^m\lambda^\ell}
 \quad(0\le k\le m).
\end{equation}
These finite data are inputs to the construction. Section~\ref{sec:parameters}
will choose them as constants depending only on the target density
$\delta$, before choosing the Label Cover soundness and alphabets.
The finite probability experiment below specifies a deterministic list
of vertices by expanding each rational probability into a multiplicity.

Choose a fixed integer $q>1/\delta$. Before invoking Label Cover,
reindex the occurring variables densely, preserving repeated occurrences
of each variable. Remove tautological clauses and duplicate literals.
A formula with no clauses maps directly to a one-vertex graph; a formula
containing an empty clause maps directly to $K_q$.
Pad a remaining short clause to width three by including all sign choices
for fresh filler variables. This preserves satisfiability at constant
overhead. In the remainder, the formula has only nonempty clauses and
has passed through this preprocessing.

\subsection{Tuples, grids, and zero-length links}

Let $\mathcal C$ be any Label Cover instance, with unused questions
removed. In layer $i\in[r]$, take all question
tuples
\[
 \mathbf q=(v_1,\ldots,v_{i-1},u_i,\ldots,u_{r-1})
 \in V^{i-1}\times U^{r-i}.
\]
The answer set of such a tuple is
\[
 M_i=\Sigma_R^{\,i-1}\times\Sigma_L^{\,r-i}.
\]
Layered products of questions with projection constraints also occur in
the multilayered PCP construction of \cite{DinurGuruswamiKhotRegev2005}.
Here each position is a separate coordinate even if question names repeat.
Different tuples in a layer use separate copies of this same answer set.

Between adjacent layers $i$ and $i+1$, impose a map whenever the
tuples agree outside position $i$ and their questions at position $i$
are the endpoints of a test occurrence $c$.
The map on answer sets applies $\pi_c$ in position $i$ and the identity
in every other position. Impose it separately for every occurrence.

Write $\Gamma=(D^{-1}\mathbb Z)/\mathbb Z$.
At each question tuple place a separate copy of $\Gamma^{M_i}$,
and let $\mathscr X$ be their disjoint union.
Make a finite undirected weighted link graph on $\mathscr X$:
\begin{itemize}
\item Within each copy, link any two points with length their sup circle
distance.
\item For every imposed map $\pi:M_i\to M_{i+1}$, link each target
point $x\in\Gamma^{M_{i+1}}$ to its source pullback $x\circ\pi$
with length zero.
\end{itemize}
Let $\dist$ be the resulting shortest-path extended pseudometric,
with distance infinity between components.
Thus distinct points may have distance zero.

Let $T$ add $1/2$ to every phase coordinate, within each copy.
Since $D$ is even, $T$ permutes $\mathscr X$ and satisfies $T^2=\id$.
It preserves all link lengths and therefore is an isometry of $\dist$.

If some $x\in\mathscr X$ satisfies
\begin{equation}\label{eq:clique-test}
 \dist(x,Tx)\le1/8,
\end{equation}
output $K_q$ and stop. This branch already has independent-set density
$1/q<\delta$. We will show that it never occurs in completeness.

\subsection{Vertices and edges}

If \eqref{eq:clique-test} fails for every point, consider the following
finite experiment.

\begin{enumerate}
\item Sample $r-1$ independent uniform occurrences
$c_1,\ldots,c_{r-1}$, with questions $(u_h,v_h)$ at position $h$.
They determine a chain of tuples
\[
 \mathbf q_i=(v_1,\ldots,v_{i-1},u_i,\ldots,u_{r-1}),
 \qquad i\in[r].
\]
Write $\pi_{ij}:M_i\to M_j$ for the corresponding composite answer maps.
They apply the test projections in positions $i,\ldots,j-1$
and leave the other positions unchanged; in particular they are
composition-consistent.

\item Independently choose $B\sim\mu$.
For every $j\in B$, independently choose $t_j$ from
\eqref{eq:coefficient-law}, and choose independent rotation entries
\[
 \widehat\theta_{j,\ell}\ \text{uniform in }\
 \{0,1/P,\ldots,(P-1)/P\},\qquad \ell\in M_j .
\]

\item With $b=\min B$, give the outcome the location
\begin{equation}\label{eq:location}
 x=\left(
 \sum_{j\in B}t_j\widehat\theta_{j,\pi_{bj}(\kappa)}
        \pmod1\right)_{\kappa\in M_b}
 \in\Gamma^{M_b}
\end{equation}
in the copy at tuple $\mathbf q_b$.
\end{enumerate}

Take distinct output vertices for these elementary outcomes with
integer multiplicities that make the experiment their exact uniform law.
The existence and size of these multiplicities are verified below.
Different vertices may have the same location.
For distinct vertices with locations $x,y$, put an edge precisely when
\begin{equation}\label{eq:edge}
 \dist(x,Ty)\le1/8.
\end{equation}
This is symmetric because $T$ is an involutive isometry:
$\dist(x,Ty)=\dist(Tx,y)=\dist(y,Tx)$.

\begin{proposition}[Explicit deterministic encoding]\label{prop:encoding}
For fixed choices of the parameters and Label Cover alphabets, this rule
constructs a nonempty finite simple
unweighted graph in deterministic polynomial time and output size,
measured in the size of the explicit Label Cover instance. Composing
with Theorem~\ref{thm:label-cover} for any fixed $\sigma$ gives polynomial
time and output size in the binary encoding length of the original
$3$-CNF formula.
\end{proposition}

\begin{proof}
Let $N=|\mathcal C|\ge1$.
All alphabets, layer counts, and grids are fixed constants.
There are polynomially many question tuples, each with a constant-size
grid. The imposed maps and links are therefore polynomially enumerable.
Every finite link length is an integer multiple of $1/D$.
After multiplying lengths by $D$, exact integer shortest-path algorithms
compute all distances, with infinity stored separately.
Shortest paths have polynomial bit length, and comparison with $1/8$
is exact by multiplication by $8$. This also implements the clique test.

There are exactly $N^{r-1}$ equally likely occurrence chains.
For each chain, the remaining outcome probabilities are fixed rational
numbers. The number of rotation entries can vary with $B$ and its layers,
but it ranges over fixed constants. Choose a common positive denominator
$C_0$ for all these conditional probabilities. For a conditional outcome
of probability $p$, emit $C_0p$ distinct vertices, omitting zero-probability
outcomes. Every chain then contributes exactly $C_0$ vertices.
The total is $C_0N^{r-1}$, and uniform sampling of these vertices is
exactly the stated experiment, followed by a uniform choice among the
copies of the chosen outcome. Exhaustive enumeration uses no random
bits. Evaluating \eqref{eq:edge} for every distinct pair remains
polynomial and gives an explicit edge list with no loops or multiple edges.

The Label Cover reduction is polynomial for the fixed $\sigma$.
The initial variable reindexing and clause preprocessing take polynomial
time in the ordinary binary input length; the two trivial branches have
constant-size outputs. Thus the construction covers every explicitly
encoded $3$-CNF formula.
\end{proof}

\begin{lemma}[All-vertex completeness]\label{lem:completeness}
If the Label Cover instance has value $1$, the clique branch does not
occur and the constructed graph has a proper three-coloring.
\end{lemma}

\begin{proof}
Fix a labeling satisfying every occurrence. At each question tuple,
form the product of its assigned answers, an element of that tuple's
answer set. Evaluate a grid point at this product answer.
This defines a map $\varphi:\mathscr X\to\T$.

Within a copy, evaluation is nonexpanding for the sup circle distance.
Across each zero link the two evaluations agree because the labeling
satisfies that test occurrence. Hence $\varphi$ is nonexpanding for
every path, and thus for $\dist$.
Also $\varphi(Tx)=\varphi(x)+1/2$.
Consequently $\dist(x,Tx)\ge1/2$ for every $x$, so the clique test fails.

If output vertices at $x,y$ are adjacent, then
\[
 d_{\T}\bigl(\varphi(x),\varphi(y)+1/2\bigr)\le1/8,
 \qquad
 d_{\T}\bigl(\varphi(x),\varphi(y)\bigr)\ge3/8>1/3.
\]
Color each vertex according to which of the half-open intervals
$[0,1/3)$, $[1/3,2/3)$, $[2/3,1)$ contains its phase.
Two phases in the same interval have circle distance less than $1/3$,
so this colors every vertex properly. All multiplicity copies receive
the color of their location.
\end{proof}

Figure~\ref{fig:geometry} shows how the two parts of the construction
serve this coloring: projection links preserve the evaluated phase,
whereas an output edge separates its endpoint phases by more than the
length of a color interval.

\begin{figure}[t]
\centering
\begin{tikzpicture}[>=Stealth,every node/.style={font=\small}]
\node (source) at (0,1.1) {$x\circ\pi\in\Gamma^{M_i}$};
\node (target) at (4.0,1.1) {$x\in\Gamma^{M_{i+1}}$};
\draw[<-] (source) -- node[above] {pullback} node[below] {zero-length link} (target);
\node (seval) at (0,-0.4) {$x(\pi(\kappa_i))$};
\node (teval) at (4.0,-0.4) {$x(\kappa_{i+1})$};
\draw[->] (source) -- node[left] {evaluate} (seval);
\draw[->] (target) -- node[right] {evaluate} (teval);
\draw (seval) -- node[above] {$=$} (teval);
\node at (2.0,-1.1) {$\pi(\kappa_i)=\kappa_{i+1}$};
\begin{scope}[shift={(8.0,0)},scale=1.15]
\draw[line width=3pt,blue!65!black] (1,0) arc[start angle=0,end angle=120,radius=1];
\draw[line width=3pt,green!45!black] (120:1) arc[start angle=120,end angle=240,radius=1];
\draw[line width=3pt,red!65!black] (240:1) arc[start angle=240,end angle=360,radius=1];
\draw[densely dashed] (30:1) -- (210:1);
\fill (30:1) circle[radius=1.5pt];
\fill (210:1) circle[radius=1.5pt];
\node[right] at (30:1.05) {$\varphi(x)$};
\node[below left] at (210:1.05) {$\varphi(x)+\tfrac12$};
\draw[line width=1.3pt] (165:1.18) arc[start angle=165,end angle=255,radius=1.18];
\node at (8:1.25) {$0$};
\node[above left] at (120:1.07) {$\tfrac13$};
\node[below left] at (240:1.07) {$\tfrac23$};
\node[align=center] at (0,-1.85) {$d_{\T}(\varphi(x),\varphi(y))\ge\tfrac38$\\for every output edge $xy$};
\end{scope}
\end{tikzpicture}
\caption{The auxiliary geometry and the output coloring. Left: a satisfying
product answer makes evaluation agree across a projection link. Here
$\kappa_i$ and $\kappa_{i+1}$ are the satisfying product answers at the
two tuples. The answer
map $\pi:M_i\to M_{i+1}$ and its phase pullback go in opposite directions.
Right: the three arcs represent the half-open color intervals. For the
illustrated phase $\varphi(x)$, the black outer arc contains the possible
phases $\varphi(y)$ of its neighbors, within $1/8$ of the antipodal phase.
The weighted links on the left define the pseudometric; the output graph
uses the separate half-shift edge rule.}
\label{fig:geometry}
\end{figure}
